\section{Introduction}
\label{sec:introduction}

Class imbalance arises when one class is substantially less frequent than the
other, often degrading a model's ability to recognize the minority class. Under
severe skew, a classifier may achieve deceptively high overall accuracy while
exhibiting poor minority-class recall, because majority predictions dominate the
evaluation.

To address this, practitioners commonly turn to three kinds of correction.
Data-level resampling techniques, such as SMOTE~\cite{chawla2002smote} and random
oversampling, add or duplicate minority examples before training begins.
Cost-sensitive reweighting~\cite{elkan2001costsensitive} instead leaves the data
untouched and penalizes mistakes on the minority class more heavily during
training. A third kind, threshold and posterior-probability correction such as
DistPFN~\cite{lee2026distpfn}, is suited to tabular foundation models such as
TabPFN~\cite{hollmann2023tabpfn,hollmann2025tabpfnv2}: instead of touching the data
or the training process, it adjusts the model's predictions after they have
already been made.

Many imbalance-correction strategies are evaluated under an implicit condition:
that the labels used to identify and characterize the minority class are
sufficiently reliable. Although class imbalance and label noise have each been
studied extensively, considerably less is known about whether standard imbalance
corrections remain beneficial when minority-label information is systematically
corrupted (Section~\ref{sec:related-work}).

If that assumption does not hold, none of these corrections can tell a genuine
minority example from a mislabeled one. A resampling method will duplicate or
synthesize around a mislabeled row exactly as readily as a correctly labeled one. A
cost-sensitive method will derive its penalty from the observed class ratio, so any
corruption in that ratio carries directly into the penalty. A posterior-rescaling
method will correct a prediction against an observed class prior that carries the
same corruption. In every case, the correction can end up amplifying the
mislabeled examples it was built to compensate for, rather than correcting for
genuine class imbalance.

This paper tests that failure mode directly, measuring when standard
imbalance-correction methods help and when they backfire under controlled, directed
label contamination. We organize this investigation around four research
questions.

\begin{itemize}
\item \textbf{RQ1.} Do standard imbalance corrections remain beneficial once
directed label error is present?
\item \textbf{RQ2.} Does the direction of label corruption change how often a
correction fails?
\item \textbf{RQ3.} At what contamination rate does each correction stop paying off?
\item \textbf{RQ4.} How much of a correction's failure comes from using a
corrupted prior instead of the true, oracle prior?
\end{itemize}

Concretely, we audit seven correction methods across two widely deployed tabular
model families, TabPFN and XGBoost~\cite{chen2016xgboost}, under two directed,
asymmetric label-error mechanisms and a range of imbalance levels, using ten real
datasets drawn from the UCI Machine Learning Repository~\cite{kelly_uci}. The design
holds the total amount of training evidence fixed across imbalance levels, so that
imbalance and label noise can be varied independently rather than confounded with
one another. Section~\ref{sec:methods} gives the full experimental setup, including
the exact corruption procedure and every configuration tested.

To answer these questions, we judge whether a correction helps or hurts, rather
than only how accurate it is, using four purpose-built metrics. Correction Gain and
Correction Harm Rate answer RQ1 and RQ2, measuring whether and how often a
correction underperforms its uncorrected baseline. The break-even contamination
rate answers RQ3, and the Prior Sensitivity Gap answers RQ4. Each metric is
defined precisely in Section~\ref{sec:methods}.

Our central finding is that every correction method tested underperforms its
matched uncorrected baseline in at least some conditions. How often this happens
depends strongly on both the direction of the corruption and the dataset the
correction is applied to, and no correction remains non-degrading across every
tested condition in both corruption directions. The complete results, including
per-method failure rates and break-even points, are reported in
Section~\ref{sec:results}.

Our contributions are as follows.

\begin{itemize}
\item We introduce a controlled experimental design that holds total training-set
size fixed across imbalance ratios and injects two directed, asymmetric
label-error mechanisms on top of it, isolating the effect of label noise from the
effect of imbalance itself.
\item We audit seven correction methods across two widely used model families,
including DistPFN and DistPFN-T~\cite{lee2026distpfn}. We validate our DistPFN and
DistPFN-T implementation against the authors' published reference code and document
an implementation error identified during this verification process
(Section~\ref{sec:methods}).
\item We operationalize correction robustness using four complementary measures,
Correction Gain, Correction Harm Rate, break-even contamination rate, and Prior
Sensitivity Gap, designed to measure whether a correction remains net-beneficial
once label noise is present, rather than measuring model accuracy alone.
\item We show, across ten real datasets, that every correction method tested can
underperform doing nothing at all once contamination reaches as little as 5\%, and
that this failure is direction-dependent rather than uniform across methods or
datasets.
\end{itemize}
